\
\chapter*{Resumen}
\addcontentsline{toc}{chapter}{Resumen}
La anemia falciforme (AF) constituye uno de los trastornos monogénicos más prevalentes a nivel mundial. Su diagnóstico oportuno en entornos con recursos limitados continúa siendo un reto por los costos y la infraestructura requerida por técnicas confirmatorias como HPLC o electroforesis. En este trabajo se estudia un \emph{pipeline} de aprendizaje automático que clasifica \emph{hologramas} provenientes de microscopia digital holográfica sin lente (DLHM), evitando la reconstrucción numérica y aprovechando la textura y patrones de interferencia para distinguir eritrocitos sanos de falciformes. El enfoque integra preprocesamiento (normalización, \emph{CLAHE}), extracción de descriptores (LBP, GLCM, energía espectral vía FFT) y clasificadores lineales y no lineales, y se compara con una alternativa de aprendizaje profundo con transfer learning. Se emplea validación cruzada anidada para estimación imparcial de desempeño bajo tamaños muestrales reducidos. Adicionalmente, se exploran técnicas de interpretabilidad (p.~ej., SHAP) para analizar la contribución de rasgos. Los resultados muestran mejoras consistentes frente a baselines y demuestran la viabilidad de DLHM+ML como apoyo al cribado de AF en escenarios de alto impacto.
